% Bruchrechnung – bank/bruchrechnung/ (zusammenbau v0.4, Heft)
\einheitenkopf[t1]{Bruchrechnung}
\setcounter{aufgabe}{0}

\begin{aufgabe}{Addieren/Subtrahieren – Prüfungsaufgaben}
\begin{teile}
\teil Ein Glücksrad hat 10 gleich große Felder mit den Zahlen 1 bis 10. Die Wahrscheinlichkeit für 8, 9 oder 10 ist $\frac{3}{10}$. Wie groß ist die Wahrscheinlichkeit für eine Zahl von 1 bis 7? \hfill (P10 2014 OS)
\teil Ein Spielwürfel hat 12 Flächen mit den Zahlen 1 bis 12. Die Wahrscheinlichkeit für 1 oder 12 ist $\frac{2}{12}$. Wie groß ist die Wahrscheinlichkeit für weder 1 noch 12? Gib sie gekürzt an. \hfill (P10 2014 OS)
\end{teile}
\end{aufgabe}

\begin{aufgabe}{Dezimal multiplizieren/dividieren – Prüfungsaufgaben}
\begin{teile}
\teil Welcher Wert ist am kleinsten: $2{,}2$; $0{,}22$; $0{,}2^2$; $22\,\%$? \hfill (P10 2023 OS) \\ \leerfeld
\teil Welcher Wert ist am kleinsten: $5{,}5$; $0{,}55$; $0{,}5^2$; $55\,\%$? \hfill (P10 2023 OS) \\ \leerfeld
\end{teile}
\end{aufgabe}

\begin{aufgabe}{Punkt vor Strich – Prüfungsaufgaben}
\begin{teile}
\teil Berechne den Wert des Terms $\frac{a + b}{c}$ für $a = 4{,}5$, $b = 1{,}5$ und $c = 4$. \hfill (P10 2021 OS)
\teil Berechne den Wert des Terms $\frac{a + b}{c}$ für $a = \frac{2}{3}$, $b = \frac{1}{6}$ und $c = 5$. \hfill (P10 2021 OS)
\teil Berechne den Wert des Terms $(a + b) : c$ für $a = 7$, $b = 2{,}6$ und $c = 1{,}2$. \hfill (P10 2016 OS)
\teil Berechne den Wert des Terms $(a + b) : c$ für $a = 10$, $b = 3{,}5$ und $c = 0{,}5$. \hfill (P10 2016 OS)
\end{teile}
\end{aufgabe}
